\section*{अध्याय \ref{ch:b1:integration} --- खंड पर समाकलन}

\begin{solution}{exo:b1:integration:1}
$\dfrac{1}{x^2-4} = \dfrac{1/4}{x - 2} - \dfrac{1/4}{x+2}$ (आवरण विधि), अतः
\[
\int_0^1 \frac{\dd x}{x^2 - 4}
= \frac14\Bigl[\ln\abs{x-2} - \ln\abs{x+2}\Bigr]_0^1
= \frac14\Bigl(\ln\frac{1}{3} - \ln 1\Bigr) = -\frac{\ln 3}{4}.
\]

खंडशः ($u' = 1$, $v = \ln t$): $\int_1^{\eu} \ln t\,\dd t = \bigl[t\ln
t\bigr]_1^{\eu} - \int_1^{\eu} \dd t = \eu - (\eu - 1) = 1$।

खंडशः ($u' = \sin t$, $v = t$): $\int_0^\pi t\sin t\,\dd t = \bigl[-t\cos
t\bigr]_0^\pi + \int_0^\pi \cos t\,\dd t = \pi + 0 = \pi$।
\end{solution}

\begin{solution}{exo:b1:integration:2}
प्रतिस्थापन $u = t^2 + 1$, $\dd u = 2t\,\dd t$:
\[
\int_0^1 \frac{t\,\dd t}{(t^2+1)^2}
= \frac12 \int_1^2 \frac{\dd u}{u^2}
= \frac12\Bigl[-\frac1u\Bigr]_1^2 = \frac14 .
\]
$u = \sin t$ के साथ: $\int_0^{\pi/2} \cos^3 t\,\dd t = \int_0^{\pi/2} (1 -
\sin^2 t)\cos t\,\dd t = \bigl[\sin t - \frac{\sin^3 t}{3}\bigr]_0^{\pi/2} =
1 - \frac13 = \frac23$।
\end{solution}

\begin{solution}{exo:b1:integration:3}
$u_n = \dfrac1n \sum_{k=1}^{n} \dfrac{1}{1 + (k/n)^2}$: यह $\intcc{0}{1}$ पर
$x \mapsto \frac{1}{1+x^2}$ का \omterm{thm:b1:integration:riemann}{रीमान योग} है, अतः $u_n \to \int_0^1 \frac{\dd
x}{1+x^2} = \arctan 1 = \dfrac\pi4$।

$\ln v_n = \dfrac1n \sum_{k=1}^{n} \ln\dfrac{n+k}{n} = \dfrac 1n
\sum_{k=1}^n \ln\Bigl(1 + \dfrac kn\Bigr) \to \int_0^1 \ln(1+x)\,\dd x =
\bigl[(1+x)\ln(1+x) - x\bigr]_0^1 = 2\ln 2 - 1$। अतः $v_n \to \eu^{2\ln 2 -
1} = \dfrac 4\eu$। (पहचान की जाँच: $\frac{(2n)!}{n!\,n^n} = \prod_{k=1}^{n}
\frac{n+k}{n}$।)
\end{solution}

\begin{solution}{exo:b1:integration:4}
सीमा $0$ है। मान लीजिए $M = \sup \abs f$ और $\varepsilon \in \intoo{0}{1}$।
$1 - \varepsilon$ पर काटिए:
\[
\Bigl| \int_0^1 x^n f \Bigr|
\leq \int_0^{1 - \varepsilon} x^n \abs f + \int_{1-\varepsilon}^1
x^n \abs f
\leq M\,(1-\varepsilon)^n + M\varepsilon .
\]
चूँकि $(1 - \varepsilon)^n \to 0$ (\cref{exo:b1:seq:3}), बाएँ पक्ष का \omterm{def:b1:reals:bounds}{उच्चतम}
सीमांत प्रत्येक $\varepsilon$ के लिए $\leq M\varepsilon$ है: अतः \omterm{thm:b1:integration:def}{समाकल} $0$
की ओर जाता है।
\end{solution}

\begin{solution}{exo:b1:integration:5}
सभी $\lambda$ के लिए $Q(\lambda) = \int_a^b (f + \lambda g)^2 = \int f^2 +
2\lambda \int fg + \lambda^2 \int g^2 \geq 0$। यदि $\int g^2 = 0$, तो $g =
0$ (कठोर धनात्मकता, \cref{thm:b1:integration:props}~(4)) और असमिका $0 \leq
0$ है। अन्यथा $Q$ सच्चा द्विघात है, जो सर्वत्र $\geq 0$ है: अतः उसका
विविक्तकर $\leq 0$ है, अर्थात् $\bigl(\int fg\bigr)^2 \leq \int f^2 \int
g^2$।

समता तभी जब विविक्तकर लुप्त हो, अर्थात् तभी जब किसी $\lambda_0$ के लिए
$Q(\lambda_0) = 0$, अर्थात् $\int (f + \lambda_0 g)^2 = 0$, अर्थात् (फिर
कठोर धनात्मकता से) $f = -\lambda_0 g$: अतः समता ठीक तब है जब $f$ और $g$
अनुपाती हों।
\end{solution}

\begin{solution}{exo:b1:integration:6}
मान लीजिए $\Phi(a) = \int_a^{a+T} f$। मूल प्रमेय
(\cref{thm:b1:integration:ftc}) से $\Phi$ \omterm{def:b1:derivative:def}{अवकलनीय} है और $\Phi'(a) = f(a + T)
- f(a) = 0$: अतः अचर।

$x > 0$ के लिए $x = nT + r$, $0 \leq r < T$ ($n = \lfloor x/T \rfloor$)
लिखिए। शाल:
\[
\int_0^x f = n \int_0^T f + \int_{nT}^{nT + r} f,
\qquad
\Bigl| \int_{nT}^{nT+r} f \Bigr| \leq T \sup_{\intcc{0}{T}} \abs f
= C .
\]
फिर $\frac 1x \int_0^x f = \frac{nT}{x}\cdot\frac 1T \int_0^T f +
O\bigl(\frac 1x\bigr)$, और $\frac{nT}{x} \to 1$: अतः सीमा $\frac1T \int_0^T
f$ है।
\end{solution}

\begin{solution}{exo:b1:integration:7}
मान लीजिए $g(x) = f(x) - x$: यह \omterm{def:b1:continuity:continuous}{संतत} है, और
\[
\int_0^1 g = \int_0^1 f - \frac12 = 0 .
\]
यदि $g$ कभी लुप्त न होता, तो मध्यवर्ती मान प्रमेय उसे एक ही चिह्न पर बाध्य
कर देती (किसी \omterm{prop:b1:reals:intervals}{अंतराल} पर दोनों चिह्न लेने वाला \omterm{def:b1:continuity:continuous}{संतत} फलन लुप्त होता है); मान
लीजिए $g > 0$। तब कठोर धनात्मकता (\cref{thm:b1:integration:props}~(4), $g >
0$ पर लगाकर, $\int g > 0$ देती है): $\int g = 0$ से विरोधाभास। अतः किसी $c$
के लिए $g(c) = 0$: $f(c) = c$।
\end{solution}

\begin{solution}{exo:b1:integration:8}
\begin{enumerate}
  \item $u' = \sin t$, $v = \sin^{n-1} t$ के साथ खंडशः:
        \[
        W_n = \bigl[-\cos t\sin^{n-1}t\bigr]_0^{\pi/2}
        + (n-1)\int_0^{\pi/2} \cos^2 t\,\sin^{n-2} t\,\dd t
        = (n-1)(W_{n-2} - W_n),
        \]
        अतः $nW_n = (n-1)W_{n-2}$। $W_0 = \frac\pi2$, $W_1 = 1$ से:
        \[
        W_{2p} = \frac{(2p-1)(2p-3)\cdots 1}{(2p)(2p-2)\cdots 2}\,
        \frac{\pi}{2} = \frac{(2p)!}{4^p (p!)^2}\,\frac\pi2,
        \qquad
        W_{2p+1} = \frac{(2p)(2p-2)\cdots 2}{(2p+1)(2p-1)\cdots 3}
        = \frac{4^p (p!)^2}{(2p+1)!} .
        \]
  \item $\intoo{0}{\frac\pi2}$ पर $0 < \sin t < 1$, अतः $\sin^{n+1} <
  \sin^n$ और $(W_n)$ (कठोरतः) ह्रासमान और धनात्मक है। पुनरावृत्ति के साथ
  घेरने पर:
        \[
        \frac{n}{n+1} = \frac{W_{n+1}}{W_{n-1}}
        \leq \frac{W_{n+1}}{W_n} \leq 1
        \quad\implies\quad \frac{W_{n+1}}{W_n} \to 1 .
        \]
        अपरिवर्त्य: पुनरावृत्ति $(n+1)W_{n+1} = nW_{n-1}$ से $a_n =
        (n+1)W_{n+1}W_n$ $a_n = a_{n-1}$ संतुष्ट करता है, अतः $a_n = a_0 = 1
        \cdot W_1 W_0 = \frac\pi2$। फिर
        \[
        n W_n^2 \sim (n+1) W_{n+1} W_n = \frac\pi2
        \quad\implies\quad
        W_n \sim \sqrt{\frac{\pi}{2n}} .
        \]
\end{enumerate}
\end{solution}

\begin{solution}{exo:b1:integration:9}
\begin{enumerate}
  \item $\intoo{0}{\pi}$ पर: $x > 0$, $a - bx = b(\frac ab - x) = b(\pi - x)
  > 0$ और $\sin x > 0$, अतः समाकल्य $> 0$ है और $I_n > 0$ (कठोर धनात्मकता)।
  परिबंध: $\intcc{0}{\pi}$ पर $x \leq \pi$ और $a - bx \leq a$, अतः $P \leq
  \frac{\pi^n a^n}{n!}$ और $I_n \leq \pi\,\frac{(\pi a)^n}{n!}$, जो $0$ की
  ओर जाता है (क्रमगुणित गुणोत्तर पद को हरा देता है: वह अभिसारी चरघातांकी
  श्रेणी का व्यापक पद है, \cref{ex:b1:seq:e} देखिए); विशेष रूप से बड़े $n$
  के लिए $I_n < 1$।
  \item $x^n(a - bx)^n = \sum_{j=0}^{n} \binom nj a^{n-j} (-b)^j x^{n+j}$ का
  प्रसार कीजिए: अतः पूर्णांक $c_j$ के साथ $P = \frac{1}{n!}\sum_j c_j
  x^{n+j}$। तब $k < n$ के लिए $P^{(k)}(0) = 0$ (मूल्यांकन) और $n \leq k \leq
  2n$ के लिए $P^{(k)}(0) = \frac{k!}{n!} c_{k-n}$, जो पूर्णांक है क्योंकि
  $n! \mid k!$। इसके अतिरिक्त $P(\pi - x) = P(x)$ (प्रतिस्थापित कीजिए: $a -
  b(\pi - x) = bx$ का प्रयोग करते हुए $\pi - x$ गुणकों की अदला-बदली कर देता
  है), अतः $P^{(k)}(\pi) = \pm P^{(k)}(0)$: ये भी पूर्णांक हैं।
  \item $Q = P - P'' + P^{(4)} - \dots$ (परिमित: $P$ की घात $2n$ है) के साथ:
  $Q + Q'' = P$, और
        \[
        \bigl(Q'\sin x - Q\cos x\bigr)' = (Q + Q'')\sin x = P\sin x .
        \]
        अतः $I_n = \bigl[Q'\sin x - Q\cos x\bigr]_0^{\pi} = Q(\pi) + Q(0)$,
        जो $0$ और $\pi$ पर $P^{(2k)}$ के मानों का योग है: (2) से पूर्णांक।
  \item बड़े $n$ के लिए $I_n$ ऐसा पूर्णांक है कि $0 < I_n < 1$: असंभव। अतः
  धारणा $\pi = \frac ab$ विफल हो जाती है: $\pi$ अपरिमेय है।
\end{enumerate}
\end{solution}

\begin{solution}{exo:b1:integration:10}
$C^1$ की स्थिति: खंडशः,
\[
\int_a^b f(t)\sin\lambda t\,\dd t
= \Bigl[-f(t)\frac{\cos\lambda t}{\lambda}\Bigr]_a^b
+ \frac{1}{\lambda}\int_a^b f'(t)\cos\lambda t\,\dd t,
\]
जिसका निरपेक्ष मान $\frac{2\sup\abs f + (b - a)\sup\abs{f'}} {\lambda} \to
0$ से परिबद्ध है।

\omterm{def:b1:continuity:continuous}{संतत} स्थिति: मान लीजिए $\varepsilon > 0$ और ऐसा \omterm{def:b1:integration:step}{सोपान फलन} $\varphi$ चुनिए कि
$\abs{f - \varphi} \leq \varepsilon$ (\cref{thm:b1:integration:approx}
$\leq\varepsilon$ \omterm{prop:b1:reals:intervals}{अंतराल} के साथ $\varphi \leq f \leq \psi$ देता है;
$\varphi$ लीजिए)। तब
\[
\Bigl|\int_a^b f\sin\lambda t\Bigr|
\leq \int_a^b \abs{f - \varphi}
+ \Bigl|\int_a^b \varphi \sin\lambda t\Bigr|
\leq (b-a)\varepsilon
+ \sum_i \abs{c_i}\,\Bigl|\int_{x_{i-1}}^{x_i} \sin\lambda t\,\dd
t\Bigr| ,
\]
और प्रत्येक $\bigl|\int \sin \lambda t\bigr| = \bigl|\frac{\cos\lambda
x_{i-1} - \cos\lambda x_i}{\lambda}\bigr| \leq \frac{2}{\lambda}$: दूसरा पद
$0$ की ओर जाता है। अतः \omterm{def:b1:reals:bounds}{उच्चतम} सीमांत प्रत्येक $\varepsilon$ के लिए $\leq
(b-a)\varepsilon$ है: इसलिए सीमा $0$ है।
\end{solution}

\begin{solution}{exo:b1:integration:11}
मान लीजिए $m = \min f$ और $M = \max f$, जो चरम मान प्रमेय से प्राप्त होते
हैं। चूँकि $g \geq 0$: $m\,g \leq fg \leq M\,g$, अतः एकदिष्टता से
\[
m \int_a^b g \;\leq\; \int_a^b fg \;\leq\; M \int_a^b g .
\]
यदि $\int_a^b g = 0$: कठोर धनात्मकता (\cref{thm:b1:integration:props}~(4))
$g \equiv 0$ पर बाध्य कर देती है, दोनों पक्ष लुप्त हो जाते हैं, और कोई भी
$c$ काम करता है। अन्यथा $t = \frac{\int fg}{\int g}$ $\intcc{m}{M} =
f(\intcc{a}{b})$ में है
(\cref{thm:b1:continuity:evt,thm:b1:continuity:ivt}), अतः किसी $c$ के लिए $t
= f(c)$।

चिह्न महत्त्व रखता है: $f(t) = g(t) = t$ वाले $\intcc{-1}{1}$ पर $\int fg =
\int_{-1}^1 t^2 = \frac23$, जबकि प्रत्येक $c$ के लिए $f(c)\int_{-1}^1 t\,\dd
t = 0$।
\end{solution}

\begin{solution}{exo:b1:integration:12}
यदि $f \equiv 0$, तो दावा रिक्तार्थ है (हर बिंदु शून्य है)। अतः $f
\not\equiv 0$ मान लीजिए और मान लीजिए $\intoo{a}{b}$ में उसके अधिक से अधिक
$n$ भिन्न शून्य हैं; मान लीजिए $z_1 < \dots < z_m$ ($m \leq n$) वे शून्य हैं
जहाँ $f$ \emph{चिह्न बदलता है} (संभवतः कोई भी नहीं)। घात $m \leq n$ का $P(t)
= \prod_{i=1}^{m}(t - z_i)$ रखिए (रिक्त गुणनफल $= 1$)। $z_i$ से कटे प्रत्येक
उपअंतराल पर $f$ और $P$ दोनों का चिह्न अचर है, और किसी $z_i$ को पार करते समय
दोनों चिह्न पलट देते हैं: अतः गुणनफल $fP$ का पूरे $\intoo{a}{b}$ पर एक ही
अचर चिह्न है। \omterm{def:b1:continuity:continuous}{संतत} होने से, तत्समक रूप से शून्य न होने से और अचर चिह्न होने
से $\bigl|\int_a^b fP\bigr| > 0$ ($\abs{fP}$ पर कठोर धनात्मकता)। पर $\int f
P$ आघूर्णों $\int f\,t^k$, $k \leq n$ का रैखिक संयोजन है, और वे सब शून्य
हैं: विरोधाभास। अतः $f$ के $\intoo{a}{b}$ में कम से कम $n + 1$ भिन्न शून्य
हैं।
\end{solution}

\begin{solution}{pb:b1:integration:1}
\textbf{1.} मान लीजिए $N = \lceil 2c \rceil$। $n \geq N$ के लिए
$\frac{c^{n+1}/(n+1)!}{c^n/n!} = \frac{c}{n+1} \leq \frac12$, अतः $0 <
\frac{c^n}{n!} \leq \frac{c^N}{N!}\,2^{-(n - N)} \to 0$: दबाव।

\textbf{2.} $k$ नियत कीजिए; $m$ पर आगमन। $m = 0$ के लिए: $\int_0^1 x^k =
\frac{1}{k+1} = \frac{k!\,0!}{(k+1)!}$। पद, खंडशः ($u = (1-x)^m$, $v' =
x^k$):
\[
\int_0^1 x^k (1-x)^m \dd x
= \Bigl[\frac{x^{k+1}}{k+1}(1-x)^m\Bigr]_0^1
+ \frac{m}{k+1}\int_0^1 x^{k+1}(1-x)^{m-1}\dd x
= \frac{m}{k+1}\cdot\frac{(k+1)!\,(m-1)!}{(k+m+1)!} ,
\]
जो $\frac{k!\,m!}{(k+m+1)!}$ है।

\textbf{3.} $k = m = n$: $\int_0^1 (x(1-x))^n = \frac{(n!)^2}{(2n+1)!} =
\frac{1}{(2n+1)\binom{2n}{n}}$। चूँकि $\intcc{0}{1}$ पर $x(1-x) \leq
\frac14$, \omterm{thm:b1:integration:def}{समाकल} $\leq 4^{-n}$ है, जिससे $\binom{2n}{n} \geq
\frac{4^n}{2n+1}$ --- जो \cref{pb:b1:seq:1} के $\binom{2n}{n}^{1/n} \to 4$
के अनुरूप है।

\textbf{4.} समाकल्य \omterm{def:b1:continuity:continuous}{संतत} है, $\geq 0$, और $\intoo{0}{1}$ पर धनात्मक, अतः
तत्समक रूप से शून्य नहीं: इसलिए उसका \omterm{thm:b1:integration:def}{समाकल} $> 0$ है
(\cref{thm:b1:integration:props}~(4))। \omterm{def:b1:reals:bounds}{ऊपरी परिबंध}: $(x(1-x))^n \leq 4^{-n}$
और $g \leq \sup\abs g$, फिर एकदिष्टता।

\textbf{5.} $A_0 = \eu - 1$; $A_1 = [x\eu^x]_0^1 - \int_0^1 \eu^x = \eu -
(\eu - 1) = 1$। खंडशः: $A_n = [x^n \eu^x]_0^1 - n\int_0^1 x^{n-1}\eu^x = \eu
- n A_{n-1}$। परिबंध: समाकल्य धनात्मक है, अतः $A_n > 0$; और $\eu^x \leq \eu$
$A_n \leq \eu\int_0^1 x^n = \frac{\eu}{n+1}$ देता है।

\textbf{6.} $A_0 = -1 + 1\cdot\eu$। यदि $A_{n-1} = \alpha_{n-1} +
\beta_{n-1}\eu$ की प्रविष्टियाँ पूर्णांक हैं, तो
\[
A_n = \eu - n\alpha_{n-1} - n\beta_{n-1}\eu
= \underbrace{(-n\,\alpha_{n-1})}_{\alpha_n}
+ \underbrace{(1 - n\,\beta_{n-1})}_{\beta_n}\,\eu ,
\]
दोनों पूर्णांक।

\textbf{7.} यदि $\eu = \frac pq$: $q A_n = q\alpha_n + p\beta_n \in \Z$, और
बड़े $n$ के लिए $0 < qA_n \leq \frac{q\eu}{n+1} < 1$: अर्थात् $0$ और $1$ के
बीच कठोरतः स्थित पूर्णांक --- असंभव। अतः $\eu \notin \Q$।
\cref{exo:b1:seq:9} में फँसा हुआ पूर्णांक $q!\,\frac pq - q!\,a_q$ था; यहाँ
वह $qA_n$ है: वही जाल, \omterm{thm:b1:integration:def}{समाकल} ईंधन।

\textbf{8.} $n = 0$: $R_0 = \int_0^1 \eu^t = \eu - 1$, अतः $\eu = 1 + R_0$।
खंडशः ($u = \eu^t$, $v = -\frac{(1-t)^{n+1}}{n+1}$):
\[
R_n = \frac{1}{n!}\Bigl(\Bigl[-\frac{(1-t)^{n+1}}{n+1}
\eu^t\Bigr]_0^1 + \frac{1}{n+1}\int_0^1 (1-t)^{n+1}\eu^t\Bigr)
= \frac{1}{(n+1)!} + R_{n+1} ,
\]
अतः सूत्र $n$ से $n + 1$ तक आगे बढ़ता है। परिबंध: $\intcc{0}{1}$ पर $1 \leq
\eu^t \leq \eu$ और $\int_0^1 (1-t)^n = \frac{1}{n+1}$ $\frac{1}{(n+1)!} \leq
R_n \leq \frac{\eu}{(n+1)!}$ देते हैं।

\textbf{9.} $\sum_{k=0}^{9} \frac{1}{k!} = \frac{986410}{362880} =
2.71828152\dots$, और
\[
\frac{1}{10!} = 2.76\cdot10^{-7} \leq R_9 \leq
\frac{2.75}{10!} = 7.58\cdot10^{-7} ,
\]
अतः $2.7182818 \leq \eu \leq 2.7182823$: उपपत्ति सहित, छह दशमलव तक $\eu =
2.718282$ (सच्चा मान $2.7182818\dots$)।

\textbf{10.} $f(1 - x) = \frac{(1-x)^n x^n}{n!} = f(x)$। $\intoo{0}{1}$ पर:
$0 < x(1-x) \leq \frac14$, अतः $0 < f \leq \frac{1}{4^n\,n!}$।

\textbf{11.} $x^n(1-x)^n = \sum_{j=0}^{n} (-1)^j\binom nj\,x^{n+j}$, अतः
$c_j \in \Z$ के साथ $f = \frac{1}{n!}\sum_j c_j\,x^{n+j}$। इसलिए $k < n$ या
$k > 2n$ के लिए $f^{(k)}(0) = 0$, और $n \leq k \leq 2n$ के लिए: $f^{(k)}(0)
= \frac{k!}{n!}\, c_{k-n}$, जो पूर्णांक है क्योंकि $n! \mid k!$। सममिति
$f^{(k)}(1) = (-1)^k f^{(k)}(0) \in \Z$ देती है।

\textbf{12.} $b^n \pi^{2n-2k} = b^n\bigl(\frac ab\bigr)^{\,n-k} =
a^{\,n-k}\,b^{\,k} \in \Z$, अतः $G(0) = \sum_k (-1)^k a^{n-k} b^k
f^{(2k)}(0)$ और इसी प्रकार $G(1)$ प्रश्न 11 से पूर्णांक हैं।

\textbf{13.} $\pi^2 G + G''$ में $\pi^2 G$ का पद $k$ $\pi^{2n-2k+2}f^{(2k)}$
ढोता है और $G''$ का पद $j = k - 1$ $(-1)^{k-1}\pi^{2n-2k+2}f^{(2k)}$ ढोता
है: अतः पहले योग में $k = 0$ और दूसरे में $j = n$ को छोड़कर सब कुछ कट जाता
है, अर्थात्
\[
G'' + \pi^2 G = b^n\bigl(\pi^{2n+2} f + (-1)^n f^{(2n+2)}\bigr)
= b^n \pi^{2n+2} f = \pi^2 a^n f
\]
($f$ की घात $2n$ है, अतः $f^{(2n+2)} = 0$; और $b^n\pi^{2n} = a^n$)। फिर
\[
\bigl(G'\sin\pi x - \pi G\cos\pi x\bigr)'
= (G'' + \pi^2 G)\sin \pi x = \pi^2 a^n f\sin\pi x .
\]

\textbf{14.} $\intcc{0}{1}$ पर समाकलन करने पर:
\[
\pi^2 a^n \int_0^1 f\sin(\pi x)\,\dd x
= \bigl[G'\sin\pi x - \pi G\cos\pi x\bigr]_0^1
= \pi\bigl(G(1) + G(0)\bigr) ,
\]
अतः $\pi a^n \int_0^1 f\sin\pi x = G(0) + G(1) \in \Z$। $\intoo{0}{1}$ पर $f
> 0$ और $\sin\pi x > 0$: बायाँ पक्ष धनात्मक है, अतः $G(0) + G(1) \geq 1$; और
$\sin \leq 1$ प्रश्न 10 के साथ उसे $\frac{\pi a^n}{4^n\,n!}$ से परिबद्ध कर
देता है।

\textbf{15.} प्रश्न 1 से ($c = \frac a4$ के साथ) $\frac{\pi a^n}{4^n n!} \to
0$: बड़े $n$ के लिए वह $< 1$ है, जो $G(0) + G(1) \geq 1$ का विरोध करता है।
अतः कोई भिन्न $\frac ab$ $\pi^2$ के बराबर नहीं है: यही लजांद्र की प्रमेय है।
यदि $\pi$ परिमेय होता, तो $\pi^2$ भी होता: अतः $\pi \notin \Q$ --- और यह
निहितार्थ केवल इसी दिशा में चलता है ($\sqrt2$ अपरिमेय है और उसका वर्ग
परिमेय), और इसीलिए $\pi^2 \notin \Q$ \cref{exo:b1:integration:9} से कठोरतः
प्रबल है।

\textbf{16.} पूर्णांकता: प्रश्न 11--12 (सिरों के \omterm{def:b1:derivative:def}{अवकलज}); लघुता: प्रश्न 10 और
1; सेतु: प्रश्न 13--14 (दूरबीनी दोहरा \omterm{thm:b1:integration:parts}{खंडशः समाकलन})। $\frac{1}{n!}$ के बिना
सिरों के आँकड़े पूर्णांक ही रहते हैं (और भी आसानी से), पर परिबंध $\frac{\pi
a^n}{4^n}$ बन जाता है, जो $0$ की ओर तभी जाता है जब $a < 4$ --- और हर
उम्मीदवार का $a = b\,\pi^2 > 9$ है। क्रमगुणित ठीक वही है जो गुणोत्तर वृद्धि
$a^n$ से आगे निकल जाता है: क्रमगुणित नहीं, तो प्रमेय नहीं।

\textbf{17.} $a \leq 10$ के लिए पूर्णांक $G(0) + G(1)$ धनात्मक और अधिक से
अधिक $\pi\,(10/4)^n/n!$ है। $n = 7$ पर: $2.5^7 = 610.35\dots$, अतः परिबंध
$\frac{\pi \times 610.35}{5040} \approx 0.38 < 1$ है ($n = 6$ पर वह अब भी
$1.07$ है): अर्थात् विरोधाभास सातवें चरण पर, स्पष्ट रूप से, आ गिरता है।

\textbf{18.} दो बार \omterm{thm:b1:integration:parts}{खंडशः समाकलन}:
\[
\int_0^1 x(1-x)\sin\pi x\,\dd x
= \frac1\pi\int_0^1 (1 - 2x)\cos\pi x\,\dd x
= \frac{2}{\pi^2}\int_0^1 \sin\pi x\,\dd x
= \frac{2}{\pi^2}\cdot\frac{2}{\pi} = \frac{4}{\pi^3}
\]
(कोष्ठक वाले पद लुप्त हो जाते हैं: $0, 1$ पर $x(1-x)$, और $0, 1$ पर $\sin\pi
x$)। सांकेतिक रूप से, $n = 1$ दूरबीन ($\pi$ पर कोई धारणा किए बिना) $f_1 =
x(1 - x)$, $f_1'' = -2$ के साथ $\pi^3\int_0^1 f_1\sin\pi x = -(f_1''(0) +
f_1''(1))$ कहती है: दायाँ पक्ष $4$ --- दोनों संगणनाएँ मेल खाती हैं।

\textbf{19.} $\eu^x$ $y' = y$ को हल करता है और $\sin\pi x$ $y'' = -\pi^2 y$
को: अर्थात् अचर गुणांकों वाले रैखिक समीकरण, अतः \omterm{thm:b1:integration:parts}{खंडशः समाकलन} अष्टक को उसी पर
लौटा देता है और सारे सीमांत आँकड़े $\Z + \Z\eu$ (क्रमशः $\pi^2$ में पूर्णांक
बहुपदों) के भीतर रखता है। मशीन को यही संवृतता चाहिए। कई बिंदुओं के लिए एक
साथ ढाले गए अष्टकों से परिष्कृत होकर वही क्रियाविधि एरमीट की प्रमेय ($\eu$
अबीजीय, 1873) और लिंडेमान की प्रमेय ($\pi$ अबीजीय, 1882) देती है --- जो इस
खंड से परे हैं।

\textbf{20.} \omterm{def:b1:poly:def}{बहुपद} भाग (अथवा गुणा करके जाँच लीजिए):
\[
x^4(1-x)^4 = (x^6 - 4x^5 + 5x^4 - 4x^2 + 4)(1 + x^2) - 4 .
\]
दिखाई गई सर्वसमिका को $1 + x^2$ से भाग देकर समाकलन करने पर:
\[
\int_0^1 \frac{x^4(1-x)^4}{1+x^2}\dd x
= \Bigl(\frac17 - \frac46 + 1 - \frac43 + 4\Bigr)
- 4\arctan 1
= \frac{22}{7} - \pi ,
\]
जहाँ $\frac17 - \frac23 + 1 - \frac43 + 4 = 3 + \frac17$ और $\arctan 1 =
\frac\pi4$ प्रयुक्त हुए।

\textbf{21.} समाकल्य \omterm{def:b1:continuity:continuous}{संतत} है और $\intoo{0}{1}$ पर धनात्मक: अतः \omterm{thm:b1:integration:def}{समाकल} $> 0$
है, इसलिए $\pi < \frac{22}{7}$। इसके अतिरिक्त $\intcc{0}{1}$ पर $\frac12
\leq \frac{1}{1+x^2} \leq 1$ और $\int_0^1 (x(1-x))^4 = \frac{(4!)^2}{9!} =
\frac{1}{630}$ (प्रश्न 3):
\[
\frac{1}{1260} \leq \frac{22}{7} - \pi \leq \frac{1}{630}
\quad\Longrightarrow\quad
3.14126 < \frac{22}{7} - \frac{1}{630} \leq \pi \leq
\frac{22}{7} - \frac{1}{1260} < 3.14207 .
\]

\textbf{22.} $x^2 + 1$ के सापेक्ष: $x^2 \equiv -1$, अतः $x^{4m} = (x^2)^{2m}
\equiv 1$ और $(1 - x)^2 = 1 - 2x + x^2 \equiv -2x$, इसलिए $(1-x)^{4m} \equiv
(-2x)^{2m} = 4^m (x^2)^m \equiv (-4)^m$: अर्थात् शेषफल अचर $(-4)^m$ है, और
भागफल $Q_m$ के गुणांक पूर्णांक हैं (किसी \omterm{def:b1:poly:def}{इकाई-अग्र} पूर्णांक \omterm{def:b1:poly:def}{बहुपद} से भाग)।
सर्वसमिका को $1 + x^2$ से भाग देकर समाकलन करने पर:
\[
J_m := \int_0^1 \frac{(x(1-x))^{4m}}{1+x^2}\dd x
= s_m + (-4)^m\,\frac{\pi}{4},
\qquad s_m = \int_0^1 Q_m \in \Q .
\]
$\pi$ के लिए हल करने पर: $r_m = (-1)^{m+1}\,4^{\,1-m} s_m \in \Q$ के साथ
$\;\abs{\pi - r_m} = 4^{\,1-m} J_m \leq 4^{\,1-m}\cdot4^{-4m} = 4^{\,1-5m}$।
$m = 1$ के लिए: $s_1 = \frac{22}{7}$, $r_1 = \frac{22}{7}$, परिबंध $4^{-4} =
\frac{1}{256}$ --- अर्थात् फिर प्रश्न 20--21।

\textbf{23.} $\bigl|\pi - \frac{22}{7}\bigr| = \frac{22}{7} - \pi \approx
1.26\cdot10^{-3}$, जो डिरिक्ले (\cref{pb:b1:derivative:1}, प्रश्न 4) की
गारंटी वाले कोटि-$2$ मानदंड $\frac{1}{7^2} \approx 2.0\cdot10^{-2}$ से सोलह
गुना बेहतर है; और $\bigl|\pi - \frac{355}{113}\bigr| \approx
2.7\cdot10^{-7}$ $\frac{1}{113^2} \approx 7.8\cdot10^{-5}$ को गुणक $\approx
300$ से हरा देता है। सिद्ध की गई किसी बात से कोई विरोध नहीं: लिउविल असमिकाएँ
सन्निकटन-त्रुटियों को केवल बीजीय संख्याओं के लिए \emph{नीचे से परिबद्ध} करती
हैं, और $\pi$ के लिए इस स्तर पर ऐसा कोई परिबंध उपलब्ध नहीं है --- अतः $\pi$
शानदार ढंग से सन्निकटित होने के लिए स्वतंत्र है।

\textbf{24.} प्रमेयिका: मान लीजिए $x = \frac pq$ और $0 < \abs{a_n + b_n x}
\to 0$। तब $\abs{a_n + b_n x} = \frac{\abs{q a_n + p b_n}}{q}$, जहाँ $q a_n
+ p b_n$ अशून्य पूर्णांक है (अशून्य इसलिए कि निरपेक्ष मान $> 0$ है): अतः
प्रत्येक $n$ के लिए $\abs{a_n + b_n x} \geq \frac1q$, जो $0$ की ओर अभिसरण का
विरोध करता है। उदाहरण: प्रश्न 7 ($x = \eu$, $a_n = \alpha_n$, $b_n =
\beta_n$); \cref{exo:b1:seq:9} (फिर $x = \eu$, $a_q = -q!\sum_{k \leq
q}\frac{1}{k!}$, $b_q = q!$ के साथ); और \cref{pb:b1:derivative:1}, प्रश्न 1,
उसका गुणोत्तर रूप है। भाग III में और \cref{exo:b1:integration:9} में यह जाल
विरोधाभास के \emph{भीतर} चलता है: परिमेयता मान लेने पर व्यंजक पूर्णांक बन
जाता है, जिसे फिर विश्लेषण $\intoo{0}{1}$ में दबा देता है --- वही सिद्धांत,
बस स्थानांतरित।

\textbf{25.} (क) पूर्णांकता बहुपदों के सिरों के कलन में रहती है (प्रश्न 6,
11--12), लघुता उन परिबंधों में जो $\sup$-एकदिष्टता देती है (प्रश्न 4, 10),
और सेतु \omterm{thm:b1:integration:parts}{खंडशः समाकलन} में (प्रश्न 8, 13--14) --- तीनों इसी अध्याय की प्रमेयें
हैं। (ख) \omterm{thm:b1:integration:def}{समाकल} वह देता है जो माध्य मान प्रमेय नहीं दे सकती थी: विश्लेषणात्मक
वस्तु और अंकगणितीय आँकड़ों के बीच एक \emph{यथार्थ} सर्वसमिका (समता, न कि
किसी अज्ञात $c$ वाली मात्र असमिका), और इसीलिए यह मशीन $\pi^2$ तक पहुँच जाती
है जबकि \cref{pb:b1:derivative:1} केवल सन्निकटन-घातांकों तक पहुँचा था। (ग)
निकाले गए परिणाम: $\eu \notin \Q$, $\pi^2 \notin \Q$ (अतः $\pi \notin \Q$),
प्रमाणित $\eu = 2.718282$, $\frac{22}{7} - \frac{1}{630} \leq \pi \leq
\frac{22}{7} - \frac{1}{1260}$, और $\binom{2n}{n} \geq \frac{4^n}{2n+1}$।
(घ) सीमांत: एरमीट और लिंडेमान इसी मशीन को अबीजीयता तक ले जाते हैं; और आपेरी
(1979) ने पूर्णांक-जाल $\zeta(3)$ पर चलाया --- यह मशीन आज भी बीसवीं शताब्दी
का गणित उत्पन्न कर रही है।
\end{solution}
